\section{Chapitre~\ref{sec:sleep}. Le sommeil}

Les modes de sommeil, les répétitions et la diminution des synapses sont mesurés par des enregistrements de neurones, par microdialyse et par microscopie électronique; l'horaire du sommeil et la bascule lente sont décrits par les modèles du livre, tandis que la fonction du sommeil relève surtout d'hypothèses.

{\footnotesize
\setlength{\tabcolsep}{3pt}\renewcommand{\arraystretch}{1.15}
\begin{longtable}{>{\raggedright}p{0.5cm}>{\raggedright}p{4.0cm}>{\raggedright}p{5.6cm}>{\raggedright}p{2.3cm}>{\raggedright\arraybackslash}p{1.7cm}}
\toprule
N\textsuperscript{o} & Affirmation & Expérience et ce qui est mesuré & Source & Statut \\
\midrule\endfirsthead
\toprule
N\textsuperscript{o} & Affirmation & Expérience et ce qui est mesuré & Source & Statut \\
\midrule\endhead
1 & Pendant le sommeil, les neurones du cortex ne se taisent pas; c'est le mode de fonctionnement qui change & Longs enregistrements de neurones du cortex chez le rat: en sommeil lent, la fréquence reste non nulle, et des périodes d'activité alternent avec des périodes de silence; après une longue veille, la fréquence est plus élevée dans tous les états, après le sommeil elle est plus basse & \cite{Vyazovskiy2009} & mesuré \\
2 & \nt{ACET} est élevé le jour et en sommeil paradoxal, bas en sommeil lent (tableau~\ref{tab:sleepmodes}) & Microdialyse dans le cortex chez des chats se déplaçant librement: pendant la veille, la libération d'acétylcholine est supérieure de $100$--$175\%$ à celle du sommeil lent, et en sommeil paradoxal elle est à peu près au niveau de la veille calme & \cite{Marrosu1995} & mesuré \\
3 & \nt{NORA} et \nt{SERO} s'apaisent en sommeil lent et se taisent presque en sommeil paradoxal & Enregistrement de neurones \nt{NORA} isolés chez le rat et de neurones \nt{SERO} chez le chat selon les stades du sommeil: la fréquence est maximale pendant la veille, plus basse en sommeil lent et presque nulle en sommeil paradoxal & \cite{AstonJonesBloom1981,McGintyHarper1976} & mesuré \\
4 & En sommeil paradoxal, les muscles sont coupés par une inhibition via \nt{GABA} et \nt{GLYC} & Chez le rat, on mesurait l'activité des neurones \nt{ACET} du muscle masticateur et on leur appliquait directement des bloqueurs: la paralysie n'est levée que si l'on bloque à la fois les deux types de récepteurs \nt{GABA} et les récepteurs \nt{GLYC} & \cite{BrooksPeever2012} & mesuré \\
5 & La pression de sommeil $S$ est un condensateur à deux seuils~\eqref{eq:sleeppressure}, $\tau_r=18.2$~h, $\tau_d=4.2$~h & Modèle à deux processus; les constantes de temps sont tirées de la façon dont la puissance des ondes lentes de l'EEG croît avec la veille et décroît pendant le sommeil, et la forme des seuils de l'heure du réveil spontané & \cite{Borbely1982,Daan1984} & théorie \\
6 & La machine trouve d'elle-même $16.1$~h de veille et $8.0$~h de sommeil (fig.~\ref{fig:twoprocess}) & Calcul selon~\eqref{eq:sleeppressure} avec des seuils oscillants, script \texttt{models/sleep\_models.py}: $16.1$~h de veille et $7.9$--$8.0$~h de sommeil & --- & modèle \\
7 & Après $40$~h sans sommeil, $S=0.90$, mais le sommeil ne dure que $8.8$~h: la dette se rembourse par la profondeur, et non par la durée & Modèle: \texttt{models/sleep\_models.py} donne $S=0.90$ et $8.8$~h. Expérience: après $40.5$~h sans sommeil, chez huit personnes, la puissance des ondes lentes de l'EEG est significativement plus élevée que d'habitude lors de la première nuit de récupération & \cite{Borbely1981} & modèle \\
8 & Physiquement, $S$ est l'adénosine & Microdialyse chez le chat: l'adénosine extracellulaire, dans l'une des régions qui commandent la veille, augmente pendant la veille, augmente encore en cas de privation de sommeil et baisse pendant le sommeil. Que $S$ soit uniquement l'adénosine n'est pas démontré & \cite{PorkkaHeiskanen1997,Dunwiddie2001} & hypothèse \\
9 & En sommeil lent, le cortex bascule: quelques centaines de ms d'activité, puis le silence, moins d'une fois par seconde ($<1$~Hz, le plus souvent $0.3$--$0.4$~Hz sous anesthésie) & Enregistrement intracellulaire chez le chat anesthésié: dépolarisation de $0.8$--$1.5$~s et longue hyperpolarisation, rythme $<1$~Hz, le plus souvent $0.3$--$0.4$~Hz; la même alternance d'activité et de silence s'observe dans le sommeil naturel (ligne~1) & \cite{Steriade1993,Vyazovskiy2009} & mesuré \\
10 & La bascule est un groupe à rétroaction et fatigue~\eqref{eq:updown}; pour $b=5$, période de $1.1$~s (valeur du modèle) et activité environ $35\%$ du temps; pour $b=1$, fonctionnement régulier (fig.~\ref{fig:updown}) & Calcul, script \texttt{models/sleep\_models.py}: période $1.11$~s, fraction du temps d'activité $0.35$ (moyenne sur un nombre entier de périodes); pour $b=1$, fraction d'activité $1.0$. La période du modèle est plus courte que la période mesurée (ligne~9) & --- & modèle \\
11 & \nt{ACET} arrête la bascule et fait passer le cortex à un fonctionnement régulier & La stimulation d'un noyau de neurones \nt{ACET} chez le chat bloquait le rythme lent dans $79\%$ des neurones du cortex et les faisait passer à un fonctionnement régulier, surtout par disparition des longues hyperpolarisations. L'acétylcholine supprime les courants potassiques lents qui créent la fatigue & \cite{SteriadeAmzica1993,McCormick1992} & mesuré \\
12 & Les bouffées de rythme à $7$--$15$~Hz naissent d'une boucle \nt{GLUT}--\nt{GABA} entre le cortex et un nœud relais & Dans des tranches du nœud relais visuel du furet, les bouffées de rythme sont créées par les bouffées de réponse des cellules relais à l'inhibition venant d'un noyau \nt{GABA} voisin, qu'elles excitent elles-mêmes & \cite{vonKrosigk1993,SteriadeMcCormickSejnowski1993} & mesuré \\
13 & Les répétitions de la journée se font en accéléré: environ $20$ fois plus vite dans l'hippocampe, $6$--$7$ fois dans le cortex & En sommeil lent, les séquences répétées de neurones de l'hippocampe sont comprimées dans le temps. Après une course sur une piste, les séquences de cellules de lieu se répétaient dans le même ordre par brèves bouffées de $\sim100$~ms, comprimées environ $20$ fois; dans le cortex, compression de $6$--$7$ fois & \cite{Nadasdy1999,LeeWilson2002,Euston2007} & mesuré \\
14 & Renforcer la coordination des répétitions avec le cortex améliore la mémoire (lien causal) & Chez le rat, après l'apprentissage, une stimulation calée sur les répétitions renforçait leur couplage avec les ondes lentes et les bouffées de rythme du cortex: le lendemain, la mémoire était élevée, alors que chez les témoins elle restait au niveau du hasard & \cite{Maingret2016} & mesuré \\
15 & La fonction des répétitions est l'apprentissage entrelacé de la mémoire lente & Théorie des deux systèmes d'apprentissage et calculs sur des réseaux artificiels: l'apprentissage séquentiel abîme l'ancien, l'apprentissage entrelacé non. Il n'existe pas d'expérience directe sur le cerveau montrant que les répétitions servent précisément à cela & \cite{McClelland1995} & hypothèse \\
16 & Après le sommeil, les synapses diminuent proportionnellement à leur taille; les grosses changent à peine & Microscopie électronique tridimensionnelle de $6920$ synapses du cortex de souris: surface de contact inférieure de $\sim18\%$ après le sommeil par rapport à la veille, diminution proportionnelle à la taille, les grosses synapses stables ne changent pas & \cite{deVivo2017} & mesuré \\
17 & La tâche principale du sommeil est la renormalisation des poids & Hypothèse de l'homéostasie synaptique; les lignes 1 et 16 s'accordent avec elle, mais ne prouvent pas qu'il s'agit de la fonction principale & \cite{TononiCirelli2014} & hypothèse \\
18 & Le sommeil paradoxal est un essai sur banc du modèle du monde & Le mode (tableau~\ref{tab:sleepmodes}) est mesuré; que le réseau fasse tourner un modèle du monde est une supposition théorique, sans expérience & \cite{Hobson2009} & hypothèse \\
19 & Lavage du cerveau pendant le sommeil: question ouverte & Une expérience: pendant le sommeil, l'espace intercellulaire est $60\%$ plus large et le nettoyage plus rapide. Une autre, avec une autre méthode de mesure: pendant le sommeil et sous anesthésie, le nettoyage est nettement plus lent. Les résultats se contredisent & \cite{Xie2013,Miao2024} & hypothèse \\
20 & Sommeil local: des zones du cortex d'un animal éveillé sombrent brièvement dans le silence & Chez le rat, après une longue veille, les neurones d'une zone du cortex se taisent brièvement, comme en sommeil lent, avec une onde lente dans l'EEG local; à ces moments-là, le rat se trompe plus souvent dans la tâche & \cite{Vyazovskiy2011} & mesuré \\
\bottomrule
\end{longtable}}
